% Source: Weinberg GR, physical PDF pp. 487--492 (printed pp. 464--468).
% Physical p. 492 repeats physical p. 491 (printed p. 468) and is
% transcribed only once.
% Coverage: Chapter 14 bibliography and references 1--69.
% Figures/tables/footnotes: none.
% Status: source-reviewed and compile-clean; visually proofed.
% Uncertainties: SOURCE ERRATUM -- reference 13's printed year conflicts with
%                both the text and cited volume; it is corrected below.
%
% SOURCE ERRATUM: reference 13 prints 1927, although Section 14.6 calls this
% Hubble's 1929 discovery and volume 15 of the Proceedings is from 1929. The
% reference below uses 1929.

\chapterbackmatter{Bibliography}

Where no other references are indicated, astronomical data are usually
taken from the following work:
\par\smallskip
\noindent C.~W.~Allen, \emph{Astrophysical Quantities}, 2nd ed.\
(The Athlone Press, London, 1955).

\medskip
\noindent\textbf{Cosmology in General}\par
\begin{itemize}
  \item[\(\square\)] H.~Bondi, \emph{Cosmology}
  (Cambridge University Press, Cambridge, 1960).

  \item[\(\square\)] W.~Davidson and J.~V.~Narlikar,
  ``Cosmological Models and Their Observational Validation,'' reprinted in
  \emph{Astrophysics} (W.~A.~Benjamin, New York, 1969).

  \item[\(\square\)] O.~Heckmann and E.~Sch\"ucking,
  ``Relativistic Cosmology,'' in \emph{Gravitation: An Introduction to
  Current Research}, ed.\ by L.~Witten (Wiley, New York, 1962), p.~438.

  \item[\(\square\)] P.~W.~Hodge, \emph{Galaxies and Cosmology}
  (McGraw-Hill, New York, 1966).

  \item[\(\square\)] \emph{La Structure et l'Evolution de l'Univers},
  Eleventh Solvay Conference, ed.\ by R.~Stoops (Brussels, 1958).

  \item[\(\square\)] G.~C.~McVittie,
  \emph{General Relativity and Cosmology}
  (University of Illinois Press, Urbana, Ill., 1965).

  \item[\(\square\)] W.~Rindler, ``Relativistic Cosmology,''
  \emph{Physics Today}, November 1967, p.~23.

  \item[\(\square\)] H.~P.~Robertson and T.~W.~Noonan,
  \emph{Relativity and Cosmology}
  (W.~B.~Saunders Co., Philadelphia, 1968).

  \item[\(\square\)] E.~L.~Sch\"ucking, ``Cosmology,'' in
  \emph{Relativity Theory and Astrophysics. 1. Relativity and Cosmology},
  ed.\ by J.~Ehlers (American Mathematical Society, Providence, R.~I.,
  1967), p.~218.

  \item[\(\square\)] D.~W.~Sciama, \emph{Modern Cosmology}
  (Cambridge University Press, Cambridge, 1971).

  \item[\(\square\)] R.~C.~Tolman,
  \emph{Relativity, Thermodynamics, and Cosmology}
  (Clarendon Press, Oxford, 1934).

  \item[\(\square\)] Ya.~B.~Zeldovich, ``Survey of Modern Cosmology,'' in
  \emph{Advances in Astronomy and Astrophysics}, Vol.~3, ed.\ by Z.~Kopal
  (Academic Press, New York, 1965).
\end{itemize}

\medskip
\noindent\textbf{History of Astronomy and Cosmology in the Twentieth
Century}\par
\begin{itemize}
  \item[\(\square\)] W.~Baade, \emph{Evolution of Stars and Galaxies},
  ed.\ by C.~Payne-Gaposchkin
  (Harvard University Press, Cambridge, Mass., 1963).

  \item[\(\square\)] F.~P.~Dickson, \emph{The Bowl of Night}
  (M.I.T.\ Press, Cambridge, Mass., 1968).

  \item[\(\square\)] J.~D.~Fernie,
  ``The Period-Luminosity Relation: A Historical Review,''
  \emph{Pub.\ Astron.\ Soc.\ Pac.}\ \textbf{81}, 707 (1969).

  \item[\(\square\)] J.~D.~North, \emph{The Measure of the Universe}
  (Clarendon Press, Oxford, 1965).

  \item[\(\square\)] \emph{Source Book in Astronomy 1900--1950},
  ed.\ by H.~Shapley (Harvard University Press, Cambridge, Mass., 1960).
\end{itemize}

\medskip
\noindent\textbf{The Cosmic Distance Scale and the Hubble Program}\par
\begin{itemize}
  \item[\(\square\)] \emph{Basic Astronomical Data},
  ed.\ by K.~Aa.~Strand (University of Chicago Press, Chicago, 1963).

  \item[\(\square\)] A.~Sandage,
  ``Cosmology---A Search for Two Numbers,''
  \emph{Physics Today}, February 1970, p.~34.

  \item[\(\square\)] A.~Sandage, in
  \emph{Problems of Extragalactic Research}
  (Macmillan, New York, 1962), p.~359.

  \item[\(\square\)] A.~Sandage,
  ``The Ability of the 200-Inch Telescope to Discriminate between Selected
  World Models,'' \emph{Ap.\ J.}\ \textbf{133}, 355 (1961).
\end{itemize}

\medskip
\noindent\textbf{Quasi-Stellar Objects}\par
\begin{itemize}
  \item[\(\square\)] G.~Burbidge and M.~Burbidge,
  \emph{Quasi-Stellar Objects}
  (W.~H.~Freeman and Co., San Francisco, 1967).

  \item[\(\square\)] L.~C.~Green,
  ``Quasars Six Years Later,'' \emph{Sky and Telescope}, May 1969.

  \item[\(\square\)] M.~Schmidt,
  ``Lectures on Quasi-Stellar Objects,'' in
  \emph{Relativity Theory and Astrophysics. 1. Relativity and Cosmology},
  op.\ cit., p.~203.

  \item[\(\square\)] M.~Schmidt,
  ``Quasi-Stellar Objects,'' in
  \emph{Annual Review of Astronomy and Astrophysics}, Vol.~7,
  ed.\ by L.~Goldberg (Annual Reviews, Inc., Palo Alto, 1969), p.~527.
\end{itemize}

\medskip
\noindent\textbf{Radio Source Counts}\par
\begin{itemize}
  \item[\(\square\)] K.~Brecher, G.~Burbidge, and P.~A.~Strittmatter,
  ``Counts of Sources and Theories,''
  \emph{Comments on Astrophysics and Space Physics} \textbf{3}, 99 (1971).

  \item[\(\square\)] M.~Ryle, ``The Counts of Radio Sources,'' in
  \emph{Annual Review of Astronomy and Astrophysics}, Vol.~6,
  ed.\ by L.~Goldberg (Annual Reviews, Inc., Palo Alto, 1968), p.~249.

  \item[\(\square\)] P.~A.~G.~Scheuer,
  ``Radio Astronomy and Cosmology,'' in
  \emph{Stars and Stellar Systems, Vol.\ IX: Galaxies and the Universe},
  ed.\ by A.\ and M.~Sandage, to be published.

  \item[\(\square\)] F.~G.~Smith, ``Radio Galaxies and Quasars,'' in
  \emph{Contemporary Physics---Trieste Symposium 1968},
  ed.\ by A.~Salam, Vol.~1
  (International Atomic Energy Agency, Vienna, 1969), p.~459.
\end{itemize}

\chapterbackmatter{References}

\begin{enumerate}
  \item\phantomsection\label{ch14-ref-1}
  H.~P.~Robertson, \emph{Ap.\ J.}\ \textbf{82}, 284 (1935);
  \emph{ibid.}\ \textbf{83}, 187, 257 (1936).

  \item\phantomsection\label{ch14-ref-2}
  A.~G.~Walker, \emph{Proc.\ Lond.\ Math.\ Soc.}\ (2), \textbf{42}, 90
  (1936).

  \item\phantomsection\label{ch14-ref-3}
  A.~Friedmann, \emph{Z.\ Phys.}\ \textbf{10}, 377 (1922);
  \emph{ibid.}\ \textbf{21}, 326 (1924).

  \item\phantomsection\label{ch14-ref-4}
  K.~G\"odel, \emph{Rev.\ Mod.\ Phys.}\ \textbf{21}, 447 (1949).
  For a general classification of homogeneous anisotropic spaces, see
  A.~Taub, \emph{Ann.\ Math.}\ \textbf{53}, 474 (1951).

  \item\phantomsection\label{ch14-ref-5}
  C.~V.~I.~Charlier, \emph{Arkiv.\ Mat.\ Astr.\ Fys.}\ \textbf{4},
  No.~24 (1908); \emph{ibid.}\ \textbf{16}, No.~22 (1922).

  \item\phantomsection\label{ch14-ref-6}
  G.~de Vaucouleurs, \emph{Science} \textbf{167}, 1203 (1970).

  \item\phantomsection\label{ch14-ref-7}
  F.~Zwicky, \emph{Pub.\ Ast.\ Soc.\ Pacific} \textbf{50}, 218 (1938);
  F.~Zwicky and K.~Rudnicki, \emph{Ap.\ J.}\ \textbf{137}, 707 (1963);
  \emph{Z.\ Astrophys.}\ \textbf{64}, 246 (1966).

  \item\phantomsection\label{ch14-ref-8}
  G.~O.~Abell, \emph{Ap.\ J.\ Suppl.}\ \textbf{3}, 211 (1958).

  \item\phantomsection\label{ch14-ref-9}
  J.~H.~Oort, XI Solvay Conference,
  \emph{La Structure et l'Evolution de l'Univers}
  (Brussels, 1958), p.~163.

  \item\phantomsection\label{ch14-ref-10}
  Table prepared for A.~S.~Eddington,
  \emph{The Mathematical Theory of Relativity}, 2nd ed.\
  (Cambridge University Press, London, 1924), p.~162.

  \item\phantomsection\label{ch14-ref-11}
  C.~Wirtz, \emph{Astr.\ Nachr.}\ \textbf{206}, 109 (1918).

  \item\phantomsection\label{ch14-ref-12}
  C.~Wirtz, \emph{Astr.\ Nachr.}\ \textbf{215}, 349 (1921);
  \emph{ibid.}\ \textbf{216}, 451 (1922);
  \emph{ibid.}\ \textbf{222}, 21 (1924);
  \emph{Scientia} \textbf{38}, 303 (1925).
  K.~Lundmark, \emph{Stock.\ Acad.\ Hand.}\ \textbf{50}, No.~8 (1920);
  \emph{Mon.\ Not.\ Roy.\ Astron.\ Soc.}\ \textbf{84}, 747 (1924);
  \emph{ibid.}\ \textbf{85}, 865 (1925).

  \item\phantomsection\label{ch14-ref-13}
  E.~P.~Hubble, \emph{Proc.\ Nat.\ Acad.\ Sci.}\ \textbf{15}, 168 (1929).

  \item\phantomsection\label{ch14-ref-14}
  K.~Schwarzschild, \emph{Vjschr.\ Astr.\ Geo.\ Lpz.}\ \textbf{35}, 337
  (1900).

  \item\phantomsection\label{ch14-ref-15}
  H.~P.~Robertson, \emph{Z.\ Astrophys.}\ \textbf{15}, 69 (1937);
  \emph{Z.\ f.\ Astrophys.}\ \textbf{15}, 69 (1938).

  \item\phantomsection\label{ch14-ref-16}
  E.~Hubble, \emph{Ap.\ J.}\ \textbf{71}, 231 (1930).

  \item\phantomsection\label{ch14-ref-17}
  See, for example, M.~Schwarzschild,
  \emph{Structure and Evolution of the Stars}
  (Princeton University Press, Princeton, N.~J., 1958), Chapter~IV.

  \item\phantomsection\label{ch14-ref-18}
  W.~Baade, \emph{Astron.\ J.}\ \textbf{100}, 137 (1944).

  \item\phantomsection\label{ch14-ref-19}
  P.~W.~Hodge and G.~Wallerstein,
  \emph{Pub.\ Astron.\ Soc.\ Pacific} \textbf{78}, 411 (1966).

  \item\phantomsection\label{ch14-ref-20}
  H.~Arp, \emph{Ap.\ J.}\ \textbf{135}, 311, 971 (1962);
  A.~Sandage, \emph{Ap.\ J.}\ \textbf{135}, 349 (1962);
  R.~Christy, \emph{Ap.\ J.}\ \textbf{144}, 108 (1966);
  L.~Plaut, in \emph{Galactic Structure}, ed.\ by A.~Blaauw and
  M.~Schmidt (University of Chicago Press, Chicago, Ill., 1965), p.~267.

  \item\phantomsection\label{ch14-ref-21}
  H.~S.~Leavitt, Harvard Circular No.~173 (1912); reprinted in
  \emph{Source Book of Astronomy}, ed.\ by H.~Shapley
  (Harvard University Press, Cambridge, 1966).

  \item\phantomsection\label{ch14-ref-22}
  H.~N.~Russell, \emph{Science} \textbf{37}, 651 (1913).

  \item\phantomsection\label{ch14-ref-23}
  E.~Hertzsprung, \emph{Astron.\ Nachr.}\ \textbf{196}, 201 (1913).

  \item\phantomsection\label{ch14-ref-24}
  H.~Shapley, \emph{Ap.\ J.}\ \textbf{48}, 89 (1918).

  \item\phantomsection\label{ch14-ref-25}
  R.~E.~Wilson, \emph{Ap.\ J.}\ \textbf{35}, 35 (1923);
  \emph{ibid.}\ \textbf{89}, 218 (1939).

  \item\phantomsection\label{ch14-ref-26}
  E.~P.~Hubble, \emph{Annual Reports of the Mount Wilson Observatory},
  1923--1924; \emph{Observatory} \textbf{48}, 139 (1925).

  \item\phantomsection\label{ch14-ref-27}
  W.~Baade, \emph{Trans.\ Int.\ Astron.\ Un.}\ \textbf{8}, 397 (1952).

  \item\phantomsection\label{ch14-ref-28}
  J.~D.~Fernie, \emph{Pub.\ Ast.\ Soc.\ Pac.}\ \textbf{81}, 707 (1969).

  \item\phantomsection\label{ch14-ref-29}
  G.~Wallerstein, \emph{Ap.\ J.}\ \textbf{127}, 583 (1958).

  \item\phantomsection\label{ch14-ref-30}
  R.~P.~Kraft, \emph{Ap.\ J.}\ \textbf{134}, 616 (1961).
  Also see R.~P.~Kraft and M.~Schmidt,
  \emph{Ap.\ J.}\ \textbf{137}, 249 (1962);
  J.~D.~Fernie, \emph{Astron.\ J.}\ \textbf{72}, 1327 (1967).

  \item\phantomsection\label{ch14-ref-31}
  A.~Sandage and G.~A.~Tammann,
  \emph{Ap.\ J.}\ \textbf{151}, 531 (1968).
  A revision of this calibration has been suggested by J.~Jung,
  \emph{Astron.\ and Astrophys.}\ \textbf{6}, 130 (1970).
  Also see A.~Sandage and G.~A.~Tammann,
  \emph{Ap.\ J.}\ \textbf{167}, 293 (1971).

  \item\phantomsection\label{ch14-ref-32}
  For a summary, see P.~W.~Hodge, \emph{Galaxies and Cosmology}
  (McGraw-Hill, New York, 1966), Chapter~12.

  \item\phantomsection\label{ch14-ref-33}
  E.~Hubble, \emph{Ap.\ J.}\ \textbf{84}, 270 (1936).

  \item\phantomsection\label{ch14-ref-34}
  A.~Sandage, \emph{Ap.\ J.}\ \textbf{152}, L149 (1968);
  \emph{Observatory} \textbf{88}, 91 (1968).

  \item\phantomsection\label{ch14-ref-35}
  R.~Racine, \emph{J.R.A.S.\ (Canada)}, in press.

  \item\phantomsection\label{ch14-ref-36}
  G.~de Vaucouleurs, \emph{Ap.\ J.}\ \textbf{159}, 435 (1970).

  \item\phantomsection\label{ch14-ref-37}
  E.~L.~Scott, \emph{Ap.\ J.}\ \textbf{62}, 248 (1957).

  \item\phantomsection\label{ch14-ref-38}
  P.~J.~E.~Peebles, \emph{Ap.\ J.}\ \textbf{153}, 13 (1968);
  J.~V.~Peach, \emph{Nature} \textbf{223}, 1140 (1969);
  P.~J.~E.~Peebles, \emph{Nature} \textbf{224}, 1093 (1969);
  B.~A.~Peterson, \emph{Ap.\ J.}\ \textbf{159}, 333 (1970);
  B.~A.~Peterson, \emph{Nature} \textbf{227}, 54 (1970); and so on.

  \item\phantomsection\label{ch14-ref-39}
  J.~B.~Oke and A.~Sandage, \emph{Ap.\ J.}\ \textbf{154}, 21 (1968).

  \item\phantomsection\label{ch14-ref-40}
  W.~A.~Baum, \emph{Ap.\ J.}\ \textbf{62}, 6 (1957).

  \item\phantomsection\label{ch14-ref-41}
  A.~Sandage, \emph{Ap.\ J.}\ \textbf{152}, L149 (1968).

  \item\phantomsection\label{ch14-ref-42}
  J.~N.~Bahcall and R.~M.~May, \emph{Ap.\ J.}\ \textbf{152}, 89 (1968).

  \item\phantomsection\label{ch14-ref-43}
  G.~de Vaucouleurs, \emph{Ap.\ J.}\ \textbf{63}, 253 (1968);
  J.~Kristian and R.~K.~Sachs,
  \emph{Ap.\ J.}\ \textbf{143}, 379 (1966).

  \item\phantomsection\label{ch14-ref-44}
  A.~Sandage, \emph{Physics Today}, February 1970, p.~34.

  \item\phantomsection\label{ch14-ref-45}
  A.~Sandage, \emph{Observatory} \textbf{88}, 91 (1968).

  \item\phantomsection\label{ch14-ref-46}
  E.~P.~Hubble and M.~L.~Humason,
  \emph{Ap.\ J.}\ \textbf{74}, 43 (1931).

  \item\phantomsection\label{ch14-ref-47}
  M.~L.~Humason, \emph{Ap.\ J.}\ \textbf{83}, 10 (1936).

  \item\phantomsection\label{ch14-ref-48}
  E.~P.~Hubble, \emph{Ap.\ J.}\ \textbf{84}, 158, 270, 516 (1936).

  \item\phantomsection\label{ch14-ref-49}
  M.~L.~Humason, N.~U.~Mayall, and A.~R.~Sandage,
  \emph{Astron.\ J.}\ \textbf{61}, 97 (1956).

  \item\phantomsection\label{ch14-ref-50}
  A.~Sandage, \emph{Ap.\ J.}\ \textbf{127}, 513 (1958).

  \item\phantomsection\label{ch14-ref-51}
  A.~Sandage, \emph{I.A.U.\ Symposium No.\ 15}, 359 (1961).

  \item\phantomsection\label{ch14-ref-52}
  A.~Sandage, \emph{Ap.\ J.}\ \textbf{134}, 916 (1961).

  \item\phantomsection\label{ch14-ref-53}
  R.~Minkowski, \emph{Ap.\ J.}\ \textbf{132}, 908 (1960).

  \item\phantomsection\label{ch14-ref-54}
  J.~V.~Peach, \emph{Ap.\ J.}\ \textbf{159}, 753 (1970).

  \item\phantomsection\label{ch14-ref-55}
  A.~Sandage, \emph{Yearbook of the Carnegie Institute of Washington}
  \textbf{65}, 163 (1966).

  \item\phantomsection\label{ch14-ref-56}
  M.~Schmidt, \emph{Nature} \textbf{197}, 1040 (1963).

  \item\phantomsection\label{ch14-ref-57}
  For a review, see G.~Burbidge and M.~Burbidge,
  \emph{Quasi-Stellar Objects}
  (W.~H.~Freeman and Co., San Francisco, 1967).

  \item\phantomsection\label{ch14-ref-58}
  F.~Hoyle and G.~R.~Burbidge, \emph{Nature} \textbf{210}, 1346 (1966).

  \item\phantomsection\label{ch14-ref-59}
  G.~R.~Burbidge and M.~Burbidge,
  \emph{Ap.\ J.}\ \textbf{148}, L107 (1967).

  \item\phantomsection\label{ch14-ref-60}
  S.~Weinberg, \emph{Ap.\ J.}\ \textbf{161}, L233 (1970).

  \item\phantomsection\label{ch14-ref-61}
  A.~Sandage, \emph{Ap.\ J.}\ \textbf{136}, 319 (1962).

  \item\phantomsection\label{ch14-ref-62}
  K.~I.~Kellermann, \emph{Ap.\ J.}\ \textbf{140}, 969 (1964);
  I.~I.~Pauliny-Toth, C.~M.~Wade, and D.~S.~Heeschen,
  \emph{Ap.\ J.\ Suppl.}\ \textbf{13}, 65 (1966).

  \item\phantomsection\label{ch14-ref-63}
  J.~F.~R.~Gower, \emph{Mon.\ Not.\ Roy.\ Astron.\ Soc.}\ \textbf{133},
  151 (1966);
  M.~Ryle, \emph{Ann.\ Rev.\ Astron.\ and Astrophys.}\ \textbf{6}, 249
  (1968).
  There are source counts at higher frequency that seem to agree with an
  \(S^{-3/2}\) law; see K.~I.~Kellermann, M.~M.~Davis, and
  I.~I.~Pauliny-Toth, \emph{Ap.\ J.\ Lett.}\ \textbf{170}, L1 (1971),
  and A.~T.~Shimmins, J.~Bolton, and J.~V.~Wall,
  \emph{Nature} \textbf{217}, 818 (1968).

  \item\phantomsection\label{ch14-ref-64}
  M.~S.~Longair and G.~G.~Pooley,
  \emph{Mon.\ Not.\ Roy.\ Astron.\ Soc.}\ \textbf{145}, 121 (1969).

  \item\phantomsection\label{ch14-ref-65}
  H.~Bondi and T.~Gold,
  \emph{Mon.\ Not.\ Roy.\ Astron.\ Soc.}\ \textbf{108}, 252 (1948).

  \item\phantomsection\label{ch14-ref-66}
  F.~Hoyle, \emph{Mon.\ Not.\ Roy.\ Astron.\ Soc.}\ \textbf{108}, 372
  (1948); \emph{ibid.}\ \textbf{109}, 365 (1949).

  \item\phantomsection\label{ch14-ref-67}
  M.~Schmidt, \emph{Ann.\ Rev.\ Astron.\ and Astrophys.}\ \textbf{7}, 527
  (1969); \emph{Ap.\ J.}\ \textbf{151}, 393 (1968);
  \emph{ibid.}\ \textbf{162}, 371 (1970).

  \item\phantomsection\label{ch14-ref-68}
  Evidence for the association of the quasi-stellar object PKS 2251+11
  with a galaxy of essentially equal red shift (\(z=0.323\)) is presented
  by J.~E.~Gunn, \emph{Ap.\ J.}\ \textbf{164}, L113 (1971).

  \item\phantomsection\label{ch14-ref-69}
  For a discussion of number counts of galaxies in optical astronomy, see
  P.~J.~E.~Peebles, \emph{Comments Ap.\ and Sp.\ Phys.}\ \textbf{3}, 173
  (1971).
\end{enumerate}
